% Unit 067 terjemahan Bahasa Indonesia dari chapter5.tex baris 468--577.
% Sumber: Methods of Algebra, Volume 2, commit
% 9a5803ff77dd3257484cb177f851a73770a59dd3 (CC BY 4.0).
% stable-unit-id: o014.aljabr2.chapter5.filtered-differential-spectral-sequence
% Cakupan: Bagian 5.4 lengkap.

% segment-id: o014.li2.chapter5.filtered-differential-spectral-sequence.h001
\section{Barisan spektral objek diferensial terfiltrasi}
\label{sec:filtered-diff-ss}
\hypertarget{o014.aljabr2.chapter5.filtered-differential-spectral-sequence}{ }

% segment-id: o014.li2.chapter5.filtered-differential-spectral-sequence.p001
Sepanjang bagian ini, kita bekerja dengan kategori abelian dengan translasi
$(\mathcal{A},T)$.

% segment-id: o014.li2.chapter5.filtered-differential-spectral-sequence.d001
\begin{definition}\label{def:filtered-diffop}
	\index{objek diferensial!terfiltrasi (filtered)}
	Objek $(X,d)\in\Obj((\mathcal{A},T)_d)$ yang dilengkapi filtrasi
	$\left(\mathrm{F}^\bullet X,d_{\mathrm{F}^\bullet X}\right)$ disebut
	\emph{objek diferensial terfiltrasi} pada $(\mathcal{A},T)$. Dalam keadaan
	ini, $\gr^pX$ juga dilengkapi secara alami dengan
	$\gr^pd:\gr^pX\to T\gr^pX$, sehingga
	$(\gr^pX,\gr^pd)\in\Obj((\mathcal{A},T)_d)$.
\end{definition}

% segment-id: o014.li2.chapter5.filtered-differential-spectral-sequence.p002
Karena $d_{\mathrm{F}^pX}$ pada subobjek $\mathrm{F}^pX$ sebaiknya dipahami
sebagai restriksi $d:X\to TX$, selanjutnya kita tetap menuliskannya sebagai
$d$. Definisi serupa dapat diberikan untuk filtrasi menaik, dengan pernyataan
yang sepenuhnya dual.

% segment-id: o014.li2.chapter5.filtered-differential-spectral-sequence.p003
Kita kembali ke barisan spektral. Bagian sebelumnya menjelaskan cara membangun
pasangan eksak dari objek diferensial. Sekarang perhatikan objek diferensial
terfiltrasi $(X,d,\mathrm{F}^\bullet X)$ pada $\mathcal{A}$. Tuliskan
$S:(Y^p)_p\mapsto(Y^{p+1})_p$ untuk funktor translasi standar pada
$\mathcal{A}^{\Z}$; jelas bahwa $S$ berkomutasi dengan $T$ yang bekerja pada
setiap $Y^p$. Karena itu morfisme yang akan dibahas mempunyai dua jenis
derajat: ``derajat filtrasi'' yang bersesuaian dengan $S$, dan ``derajat
internal'' yang bersesuaian dengan $T$. Untuk sementara kita berfokus pada
yang pertama.

% segment-id: o014.li2.chapter5.filtered-differential-spectral-sequence.p004
Karena filtrasi tersebut menurun, terdapat monomorfisme
\[ \alpha:\left(\mathrm{F}^{p+1}X,d\right)_{p\in\Z}
	\hookrightarrow S^{-1}\left(\mathrm{F}^{p+1}X,d\right)_{p\in\Z}
	=\left(\mathrm{F}^{p}X,d\right)_{p\in\Z}. \]
Dengan demikian,
$E_0:=\Coker(\alpha)=\left(\gr^pX,\gr^pd\right)_{p\in\Z}$ dan
$E_1^p=\Hm(\gr^pX,\gr^pd)$.

% segment-id: o014.li2.chapter5.filtered-differential-spectral-sequence.p005
Pandang $\alpha$ sebagai morfisme berderajat $-1$ dalam
$\mathcal{A}^{\Z}$ (terhadap $S$). Dengan menerapkan
Proposisi~\ref{prop:differential-exact-couple}, kita memperoleh pasangan eksak
berderajat pada $\mathcal{A}^{\Z}$:
\[\mathscr{C}=\left[\begin{tikzcd}[column sep=tiny]
	\Hm\left(\mathrm{F}^{p+1}X,d\right)_p
	\arrow[rr, "-1"', "{\Hm(\alpha)}"] & &
	\Hm\left(\mathrm{F}^{p+1}X,d\right)_p \arrow[ld, "{+1}"] \\
	& \Hm\left(\gr^pX,\gr^pd\right)_p \arrow[lu] &
\end{tikzcd}\right],\]
di mana panah $\nwarrow$ berasal dari morfisme penghubung yang diinduksi oleh
barisan eksak pendek
$0\to(\mathrm{F}^{p+1}X,d)\to(\mathrm{F}^pX,d)
\to(\gr^pX,\gr^pd)\to0$.

% segment-id: o014.li2.chapter5.filtered-differential-spectral-sequence.p006
Barisan
$\mathscr{E}=(E_r^p,d_r^p)_{\substack{r\geq0\\p\in\Z}}$ yang diperoleh
dengan cara ini disebut barisan spektral yang ditentukan oleh
$\mathrm{F}^\bullet X$ dalam $\mathcal{A}^{\Z}$. Pada akhir
\S\ref{sec:exact-couples} telah dijelaskan bahwa $d_r$ adalah morfisme
berderajat $r$ terhadap $S$, yaitu
\[ d_r=\left(d_r^p:E_r^p\to(S^rE_r)^p=E_r^{p+r}\right)_{p\in\Z}; \]
tentu saja, derajat $d_r$ juga dapat dibaca dari rumus dalam
Proposisi~\ref{prop:differential-exact-couple}.

% segment-id: o014.li2.chapter5.filtered-differential-spectral-sequence.p007
Dengan kata lain, objek diferensial terfiltrasi
$(X,d,\mathrm{F}^\bullet X)$ menghasilkan barisan spektral bergradasi
kohomologis sebagaimana dalam Definisi~\ref{def:graded-spectral-sequence}.
Perhatikan bahwa $d_r^p$ di sini mungkin juga mempunyai derajat internal
(kecuali jika $T=\identity_{\mathcal{A}}$); derajat ini tidak dicantumkan
dalam notasi dan akan dirinci dalam \S\ref{sec:filtered-cplx-ss}.

% segment-id: o014.li2.chapter5.filtered-differential-spectral-sequence.p008
Menurut pola \eqref{eqn:spectral-sequence-Z-B}, definisikan subobjek dari
$Z_0=(\gr^pX,\gr^pd)_p$ sebagai
\[ B_r=(B_r^p)_p\subset(Z_r^p)_p=Z_r. \]

% segment-id: o014.li2.chapter5.filtered-differential-spectral-sequence.p009
Untuk menyederhanakan notasi, pernyataan berikut tidak menampilkan derajat
internal $d$; ini setara dengan meninjau kasus khusus
$T=\identity_{\mathcal{A}}$. Perluasan ke kasus umum bersifat rutin: funktor
translasi cukup disisipkan pada tempat yang tepat dalam setiap kemunculan
$d^{-1}(\cdots)$ atau $d(\cdots)$ agar ungkapannya bermakna secara ketat.

% segment-id: o014.li2.chapter5.filtered-differential-spectral-sequence.t001
\begin{proposition}\label{prop:filtered-diff-E}
	Untuk objek diferensial terfiltrasi $(X,d,\mathrm{F}^\bullet X)$, barisan
	spektral terkait memenuhi
	\begin{align*}
		Z_r^p&=\dfrac{\left(\mathrm{F}^pX\cap d^{-1}\mathrm{F}^{p+r}X\right)
		+\mathrm{F}^{p+1}X}{\mathrm{F}^{p+1}X}, \\
		B_r^p&=\dfrac{\left(\mathrm{F}^pX\cap d\mathrm{F}^{p-r+1}X\right)
		+\mathrm{F}^{p+1}X}{\mathrm{F}^{p+1}X},
	\end{align*}
	dan $d_r^p:E_r^p\to E_r^{p+r}$ diinduksi oleh morfisme
	$d^{-1}\mathrm{F}^{p+r}X\xrightarrow{d}
	\mathrm{F}^{p+r}X\cap\Ker(d)$.
\end{proposition}
\begin{proof}
	Perhatikan bahwa komponen berderajat $p$ dari $\alpha^r$ dapat dipandang
	sebagai penyertaan $\mathrm{F}^pX\hookrightarrow\mathrm{F}^{p-r}X$ atau
	$\mathrm{F}^{p+r}X\hookrightarrow\mathrm{F}^pX$. Pernyataan tersebut
	merupakan versi bergradasi dari
	Proposisi~\ref{prop:differential-exact-couple}.
\end{proof}

% segment-id: o014.li2.chapter5.filtered-differential-spectral-sequence.p010
Karena $\varinjlim$ dan $\varprojlim$ dalam $\mathcal{A}^{\Z}$ diambil
derajat demi derajat, $E_\infty=(E_\infty^p)_{p\in\Z}$ dapat dituliskan
sebagai
\begin{equation}\label{eqn:filtered-diff-infty}
	E_\infty^p=\dfrac{Z_\infty^p}{B_\infty^p}
	=\dfrac{\bigcap_r\left((\mathrm{F}^pX\cap d^{-1}\mathrm{F}^{p+r}X)
	+\mathrm{F}^{p+1}X\right)}
	{\bigcup_r\left((\mathrm{F}^pX\cap d\mathrm{F}^{p-r+1}X)
	+\mathrm{F}^{p+1}X\right)},
\end{equation}
dengan syarat $\bigcap$ dan $\bigcup$ yang ditampilkan ada untuk setiap $p$.

% segment-id: o014.li2.chapter5.filtered-differential-spectral-sequence.r001
\begin{remark}
	Proposisi~\ref{prop:filtered-diff-E} mengenai barisan spektral dari objek
	diferensial terfiltrasi merupakan landasan pembahasan selanjutnya dalam bab
	ini. Proposisi tersebut merupakan penerapan pasangan eksak, tetapi rumus
	dalam Proposisi~\ref{prop:filtered-diff-E} juga dapat dipandang sebagai
	definisi langsung, lalu semua sifat yang diperlukan bagi barisan spektral
	diperiksa dari rumus itu.
\end{remark}

% segment-id: o014.li2.chapter5.filtered-differential-spectral-sequence.p011
Berdasarkan deskripsi di atas, barisan spektral $E_r^p$ dari
$(X,d,\mathrm{F}^\bullet X)$ dapat dipandang sebagai proses pendekatan bertahap
terhadap $\Hm(X,d)$. Untuk memperjelas arti ``mendekati'', kita
memerlukan konsep berikut.

% segment-id: o014.li2.chapter5.filtered-differential-spectral-sequence.d002
\begin{definition}[Filtrasi terinduksi]\label{def:induced-filtration-H}
	\index{filtrasi!terinduksi (induced)}
	Untuk suatu objek diferensial terfiltrasi
	$(X,d,\mathrm{F}^\bullet X)$ pada $(\mathcal{A},T)$, objek $\Hm(X,d)$
	dilengkapi filtrasi terinduksi
	\[ \mathrm{F}^p\Hm(X,d):=
	\Image\left[\mathrm{F}^pX\cap\Ker(d)\to\Hm(X,d)\right],
	\qquad p\in\Z. \]
\end{definition}

% segment-id: o014.li2.chapter5.filtered-differential-spectral-sequence.l001
\begin{lemma}\label{prop:induced-filtration-H}
	Jika terdapat $N$ dengan $\mathrm{F}^NX=0$, maka
	$\mathrm{F}^N\Hm(X,d)=0$. Jika $\mathrm{F}^\bullet X$ merupakan filtrasi
	menyeluruh menurut Definisi~\ref{def:filtration-properties}, maka
	$\mathrm{F}^\bullet\Hm(X,d)$ juga menyeluruh.
\end{lemma}
\begin{proof}
	Bagian pertama langsung; kita buktikan bagian kedua. Menurut
	Lema~\ref{prop:image-sum-preimage-intersection}~(ii),
	$\bigcup_p\mathrm{F}^p\Hm(X,d)$ adalah citra dari
	$\bigcup_p(\mathrm{F}^pX\cap\Ker(d))$, sedangkan
	\eqref{eqn:exhaustion-intersection} memberikan
	\[ \bigcup_p\left(\mathrm{F}^pX\cap\Ker(d)\right)
	=\left(\bigcup_p\mathrm{F}^pX\right)\cap\Ker(d)=\Ker(d). \]
	Selain itu, jika terdapat $M$ dengan $\mathrm{F}^MX=X$, tentu saja
	$\mathrm{F}^M\Hm(X,d)=\Hm(X,d)$.
\end{proof}

% segment-id: o014.li2.chapter5.filtered-differential-spectral-sequence.p012
Versi bergradasi dari pengamatan ini akan diterapkan dalam
\S\ref{sec:filtered-cplx-ss}.

% segment-id: o014.li2.chapter5.filtered-differential-spectral-sequence.l002
\begin{lemma}\label{prop:induced-filtration-gr}
	Untuk objek diferensial terfiltrasi $(X,d,\mathrm{F}^\bullet X)$ dan setiap
	$p\in\Z$, terdapat isomorfisme kanonik
	\[ \gr^p\Hm(X,d)\simeq
	\dfrac{\mathrm{F}^pX\cap\Ker(d)}
	{(\mathrm{F}^{p+1}X\cap\Ker(d))
	 +(\mathrm{F}^pX\cap\Image(T^{-1}d))}. \]
\end{lemma}
\begin{proof}
	Uraikan definisi
	$\mathrm{F}^p\Hm(X,d)/\mathrm{F}^{p+1}\Hm(X,d)$ dan gunakan teorema
	isomorfisme standar untuk kategori abelian,
	Proposisi~\ref{prop:Abel-cat-isom-thm}, untuk memperoleh
	\begin{align*}
		\gr^p\Hm(X,d)
		&=\dfrac{(\mathrm{F}^pX\cap\Ker(d))+\Image(T^{-1}d)}
		{(\mathrm{F}^{p+1}X\cap\Ker(d))+\Image(T^{-1}d)} \\
		&=\dfrac{(\mathrm{F}^pX\cap\Ker(d))
		 +(\mathrm{F}^{p+1}X\cap\Ker(d))+\Image(T^{-1}d)}
		{(\mathrm{F}^{p+1}X\cap\Ker(d))+\Image(T^{-1}d)} \\
		&\simeq\dfrac{\mathrm{F}^pX\cap\Ker(d)}
		{((\mathrm{F}^{p+1}X\cap\Ker(d))+\Image(T^{-1}d))
		 \cap(\mathrm{F}^pX\cap\Ker(d))} \\
		&=\dfrac{\mathrm{F}^pX\cap\Ker(d)}
		{(\mathrm{F}^{p+1}X\cap\Ker(d))
		 +(\mathrm{F}^pX\cap\Image(T^{-1}d))}.
	\end{align*}
	Langkah terakhir memakai fakta bahwa $\mathrm{Sub}_X$ adalah kekisi
	modular\footnote{Artinya, untuk semua subobjek $A,A^\flat,B\subset X$ dengan
	$A^\flat\subset A$, berlaku
	$A^\flat+(A\cap B)=(A^\flat+B)\cap A$.}
	(Teorema~\ref{prop:subobject-modularity}), serta
	$\mathrm{F}^{p+1}X\cap\Ker(d)\subset\mathrm{F}^pX\cap\Ker(d)$ dan
	$\Image(T^{-1}d)\subset\Ker(d)$. Langkah-langkah lainnya standar.
\end{proof}

% segment-id: o014.li2.chapter5.filtered-differential-spectral-sequence.d003
\begin{definition-proposition}\label{def:diff-obj-convergence}
	\index{barisan spektral!konvergen lemah, konvergen kuat
	(weakly convergent, strongly convergent)}
	Untuk suatu objek diferensial terfiltrasi
	$(X,d,\mathrm{F}^\bullet X)$, bangun barisan spektral bergradasi
	kohomologis terkait $(E_r,d_r)_r$. Jika $\bigcap_r$ dan $\bigcup_r$ dalam
	\eqref{eqn:filtered-diff-infty} ada, maka limit $E_\infty$ ada, dan
	$\gr\Hm(X,d)$ yang ditentukan oleh filtrasi terinduksi dapat direalisasikan
	secara kanonik sebagai subkuosien dari $E_\infty$.
	\begin{itemize}
		\item Jika $\gr\Hm(X,d)=E_\infty$, maka barisan spektral
		$(E_r,d_r)_r$ disebut \emph{konvergen lemah}.
		\item Jika $(E_r,d_r)_r$ konvergen lemah dan
		$\mathrm{F}^\bullet\Hm(X,d)$ menyeluruh serta lengkap, maka barisan
		spektral $(E_r,d_r)_r$ disebut \emph{konvergen kuat}.
	\end{itemize}
\end{definition-proposition}
\begin{proof}
	Misalkan $p\in\Z$. Dari ungkapan untuk $E_\infty^p$ dalam
	\eqref{eqn:filtered-diff-infty}, pembilangnya memuat
	$(\mathrm{F}^pX\cap\Ker(d))+\mathrm{F}^{p+1}X$, sedangkan penyebutnya
	terkandung dalam
	$(\mathrm{F}^pX\cap\Image(T^{-1}d))+\mathrm{F}^{p+1}X$
	(ingat bahwa $d$ adalah morfisme $X\to TX$). Dengan demikian diperoleh
	subkuosien berikut dari $E_\infty^p$:
	\begin{multline*}
		\dfrac{(\mathrm{F}^pX\cap\Ker(d))+\mathrm{F}^{p+1}X}
		{(\mathrm{F}^pX\cap\Image(T^{-1}d))+\mathrm{F}^{p+1}X} \\
		\simeq\dfrac{\mathrm{F}^pX\cap\Ker(d)}
		{((\mathrm{F}^pX\cap\Image(T^{-1}d))+\mathrm{F}^{p+1}X)
		 \cap(\mathrm{F}^pX\cap\Ker(d))} \\
		=\dfrac{\mathrm{F}^pX\cap\Ker(d)}
		{(\mathrm{F}^pX\cap\Image(T^{-1}d))
		 +(\mathrm{F}^{p+1}X\cap\Ker(d))}.
	\end{multline*}
	Isomorfisme pertama bersifat standar. Kesamaan berikutnya memakai fakta
	bahwa $\mathrm{Sub}_X$ adalah kekisi modular, bersama
	$\mathrm{F}^pX\cap\Image(T^{-1}d)
	\subset\mathrm{F}^pX\cap\Ker(d)$ dan
	$\mathrm{F}^{p+1}X\subset\mathrm{F}^pX$. Sekarang terapkan
	Lema~\ref{prop:induced-filtration-gr}.
\end{proof}

% segment-id: o014.li2.chapter5.filtered-differential-spectral-sequence.p013
Definisi konvergensi sedikit berbeda-beda antar-rujukan.
Semua pernyataan di atas memiliki versi yang bersesuaian untuk objek
diferensial dengan filtrasi menaik $(X,d,\mathrm{F}_\bullet X)$ dan barisan
spektral bergradasi homologis terkait.
